% ============================================================
% I. INTRODUCCIÓN
% ============================================================

\section{Introducción}

Una vez digitalizada una señal (\texttt{LAB1}), la siguiente etapa natural en una cadena de procesamiento digital es el filtrado: la supresión selectiva de bandas de frecuencia no deseadas —ruido electrónico, interferencia, componentes fuera del rango de interés físico— sin distorsionar la información útil contenida en la señal. Esta decisión de diseño involucra dos familias principales de filtros digitales convencionales, de Respuesta al Impulso Finita (FIR) y de Respuesta al Impulso Infinita (IIR), cada una con compromisos distintos entre linealidad de fase, orden/costo computacional y precisión de la respuesta en frecuencia lograda; y, para señales cuyas características estadísticas no son completamente conocidas o varían en el tiempo, familias alternativas de filtrado adaptativo y estadístico (media móvil, modelos autorregresivos, filtros adaptativos propiamente dichos).

Esta práctica aplica ambos enfoques sobre dos tipos de señales: señales periódicas generadas externamente con un generador de ondas, usadas para caracterizar el comportamiento de ganancia y fase de los filtros convencionales bajo condiciones controladas; y señales reales de sensores —una unidad inercial (IMU) y los encoders de los motorreductores— capturadas durante el movimiento de un robot Minisumo, donde el ruido y la señal de interés comparten en parte el mismo rango de frecuencias y la elección de estrategia de filtrado tiene un impacto directo en la calidad de la estimación de movimiento del robot.

El objetivo es triple: (1) diseñar, simular y verificar en línea filtros FIR e IIR que cumplan especificaciones de frecuencia de corte y atenuación dadas, caracterizando sus diferencias de desempeño en ganancia, fase y costo computacional; (2) aplicar y comparar estrategias de filtrado convencional, adaptativo y estadístico sobre señales reales de sensores, evaluando el compromiso entre ruido eliminado y fidelidad de la información conservada; y (3) integrar estas técnicas en un sistema completo —un robot Minisumo teleoperado mediante ROS2— donde el filtrado de sensores alimenta directamente la estimación de posición del robot, cerrando el ciclo entre procesamiento de señales y una tarea de ingeniería real.
